\chapter{Theoretical Background and Models}
\label{app:math_foundations}

This appendix provides expanded mathematical background and explanation for the mathematical foundations and literature review in chapter \ref{ch:literature_review}.

\section{Brownian Motion}
\label{app:sec:bm}

Brownian motion is the continuous-time stochastic process. It was first formally constructed by Wiener, which is why is also called the \textit{Wiener process}.

\textbf{Definition} A stochastic process $\{W_t\}_{t \geq 0}$ defined on a
probability space $(\Omega, \mathcal{F}, \mathbb{P})$ is called a
\textit{standard Brownian motion} if it satisfies the following four properties:
\begin{enumerate}
    \item \textit{Initialisation}: $W_0 = 0$ almost surely.
    \item \textit{Independent increments}: for any $0 \leq t_0 < t_1 < \cdots <
    t_n$, the increments $W_{t_1} - W_{t_0},\, W_{t_2} - W_{t_1},\, \ldots,\,
    W_{t_n} - W_{t_{n-1}}$ are mutually independent random variables.
    \item \textit{Stationary Gaussian increments}: for any $0 \leq s < t$,
    \[
        \label{eq:bm_increment}
        W_t - W_s \;\sim\; \mathcal{N}(0,\, t - s).
    \]
    \item \textit{Continuous paths}: the map $t \mapsto W_t$ is almost surely
    continuous.
\end{enumerate}

These four axioms are minimal and mutually consistent: any two of them imply the others \cite{oksendal2003stochastic}.

\subsection{Key Statistical Properties}

Setting $s = 0$ in property~3, we immediately obtain:
\[
    \mathbb{E}[W_t] = 0, \qquad \operatorname{Var}(W_t) = t, \qquad
    W_t \sim \mathcal{N}(0,\, t).
\]
The variance grows linearly in time: uncertainty accumulates at rate $\sqrt{t}$,
not $t$. This sub-linear growth is the signature of diffusion and contrasts with
a deterministic path (variance 0) or with a purely random variable with no time structure.

\textbf{Covariance and correlation} For $0 \leq s \leq t$, write
$W_t = W_s + (W_t - W_s)$. Since $W_t - W_s$ is independent of $W_s$ (property~2) and has mean zero: \(\operatorname{Cov}(W_s, W_t)=s=\min(s,t)\). Consequently $\operatorname{Corr}(W_s, W_t) = \sqrt{s/t}$ for $s \leq t$.
Two observations far apart in time are weakly correlated, so there is no long-range memory.

\textbf{Markov property} The future distribution of $W_t$ given all past
information $\mathcal{F}_s = \sigma(W_u : u \leq s)$ depends on the past only
through the current value $W_s$:
\[
    P(W_t \in A \mid \mathcal{F}_s) = P(W_t \in A \mid W_s), \qquad s \leq t.
\]
\textbf{Martingale property} With respect to $\{\mathcal{F}_t\}$, $W_t$ is a
martingale:
\[
    \mathbb{E}[W_t \mid \mathcal{F}_s] = W_s, \qquad s \leq t.
\]
Proof: $\mathbb{E}[W_t \mid \mathcal{F}_s] = \mathbb{E}[W_s + (W_t - W_s)
\mid \mathcal{F}_s] = W_s + \mathbb{E}[W_t - W_s] = W_s$, using independence
and zero mean of the increment, where the martingale property encodes the absence of predictable drift in the noise.

\section{It\^{o}'s Lemma and the Black-Scholes Model}
\label{app:sec:ito}

Ordinary calculus provides the chain rule: if $f$ is smooth and $x(t)$ is
differentiable, then $df = f'(x)\,dx$. Because BM paths are not differentiable,
this rule fails when $x$ is driven by a BM. It\^{o}'s Lemma provides the correct generalisation.

\subsection{Statement}

Let $\{X_t\}$ be an It\^{o} process satisfying:
\[
    dX_t = a_t\, dt + b_t\, dW_t,
\]
where $a_t$ (the drift) and $b_t$ (the diffusion coefficient) are stochastic processes satisfying standard integrability conditions
\cite{oksendal2003stochastic}. Let $f(t, x)$ be a function that is once
continuously differentiable in $t$ and twice continuously differentiable in $x$.
Then:
\[
    \label{eq:ito_lemma}
    df(t, X_t) = \left(\frac{\partial f}{\partial t}
    + a_t\, \frac{\partial f}{\partial x}
    + \frac{1}{2}\, b_t^2\,
    \frac{\partial^2 f}{\partial x^2}\right)\! dt
    + b_t\, \frac{\partial f}{\partial x}\, dW_t.
\]

The term $\frac{1}{2} b_t^2 \frac{\partial^2 f}{\partial x^2}\, dt$ is the
It\^{o} correction. It is absent in ordinary calculus and appear
because the stochastic increment has a non-negligible second-order contribution. A proof is provided in \cite{oksendal2003stochastic}, chapter 4, \emph{The Itô Formula and the Martingale Representation Theorem}

\subsection{Proof of the Black-Scholes Price Equation}

Chapter \ref{ch:literature_review} states that applying It\^{o}'s Lemma to $\log S_t$ under GBM
(Eq.~\ref{eq:blackscholes}) yields the log-normal price formula. It was derived in full here. The GBM is the It\^{o} process with drift coefficient
$a_t = \mu S_t$ and diffusion coefficient $b_t = \nu S_t$:
\[
    dS_t = \mu\, S_t\, dt + \nu\, S_t\, dW_t.
\]
Apply It\^{o}'s Lemma (Eq.~\ref{eq:ito_lemma}) to $f(S_t) = \log S_t$:
\[
    \frac{\partial f}{\partial t} = 0, \qquad
    \frac{\partial f}{\partial S} = \frac{1}{S_t}, \qquad
    \frac{\partial^2 f}{\partial S^2} = -\frac{1}{S_t^2}.
\]
Substituting $a_t = \mu S_t$ and $b_t = \nu S_t$:
\begin{align*}
    d(\log S_t)
    &= \left(0 + \mu S_t \cdot \frac{1}{S_t}
       + \frac{1}{2}(\nu S_t)^2 \cdot \left(-\frac{1}{S_t^2}\right)\right)\! dt
       + \nu S_t \cdot \frac{1}{S_t}\, dW_t \\
    &= \left(\mu - \frac{1}{2}\nu^2\right)\! dt + \nu\, dW_t.
\end{align*}
This is a Brownian motion with constant drift $\mu - \frac{1}{2}\nu^2$ and
constant diffusion coefficient $\nu$. Integrating over $[0, t]$:
\[
    \log S_t - \log S_0
    = \left(\mu - \frac{1}{2}\nu^2\right)t + \nu W_t,
\]
and exponentiating:
\[
\label{eq:gbm_solution_app}
    S_t = S_0 \exp\!\left[\left(\mu - \frac{1}{2}\nu^2\right)t
          + \nu W_t\right].
\]
Since $W_t \sim \mathcal{N}(0, t)$, the log-return over $[0, t]$ satisfies:
\[
    \label{eq:lognormal_returns_app}
    \log\!\left(\frac{S_t}{S_0}\right)
    \;\sim\; \mathcal{N}\!\left(\left(\mu - \frac{1}{2}\nu^2\right)t,\;
    \nu^2 t\right).
\]
Log-returns are normally distributed with mean proportional to $t$ and variance
proportional to $t$, confirming the statement in Section \ref{sec:classical_models_literature} and the foundation
of the Black--Scholes option pricing model \cite{black1973pricing}.


\textbf{Interpreting the It\^{o} correction} The term $-\frac{1}{2}\nu^2$ is
not an approximation error, instead it is present because the exponential map from log-price to price is convex, so randomness in log-space increases the arithmetic expectation in price-space, therefore it compensates for this convexity effect so that the price process still has drift \(\mu\) in expectation, so that \(\mathbb{E}[S_t]=S_0\, e^{\mu t}\).

\section{Poisson Process and Merton's Jump-Diffusion Model}
\label{app:sec:poisson_merton}

The Poisson process is the canonical model for events that arrive randomly,
infrequently, and with no predictable timing, these in the financial context are sudden
discontinuous jumps in asset prices.

\subsection{The Poisson Process}

\textbf{Definition} A stochastic process $\{P_t\}_{t \geq 0}$ taking values in
$\mathbb{N}_0 = \{0, 1, 2, \ldots\}$ is a \textit{Poisson process with intensity}
$\lambda > 0$ if:
\begin{enumerate}
    \item $P_0 = 0$ almost surely.
    \item Independent increments: for any $0 \leq t_0 < t_1 < \cdots < t_n$,
    the increments $P_{t_1} - P_{t_0},\, \ldots,\, P_{t_n} - P_{t_{n-1}}$ are
    mutually independent.
    \item Stationary Poisson increments: for any $0 \leq s < t$,
    \begin{equation}
        \label{eq:poisson_increment}
        P_t - P_s \;\sim\; \operatorname{Poisson}(\lambda(t - s)),
    \end{equation}
    i.e., $P(P_t - P_s = k) = \frac{e^{-\lambda(t-s)}[\lambda(t-s)]^k}{k!}$
    for $k \in \mathbb{N}_0$.
    \item At most one event per instant: $P(P_{t+h} - P_t \geq 2) = o(h)$
    as $h \to 0^+$.
\end{enumerate}

The process $P_t$ counts the number of jumps that have occurred by time $t$.
From Eq.~(\ref{eq:poisson_increment}):
\[
    \mathbb{E}[P_t] = \lambda t, \qquad \operatorname{Var}(P_t) = \lambda t.
\]
The intensity of $\lambda$ simultaneously controls the average jump rate and the variability of the jump count. Because the waiting time between jumps are independent then we are dealing with a memoryless process, where the time until the next jump does not depend on how long
one has already waited, see limitations in Sec.~\ref{app:sec:limitations}).

\subsection{Merton's Jump-Diffusion Model}

Merton \cite{merton1976option} extended the GBM (Eq.~\ref{eq:blackscholes}) by
superimposing a compound Poisson jump component onto the continuous diffusion:\(
    \label{eq:merton_sde_app}
    dS_t = \mu\, S_t\, dt + \nu\, S_t\, dW_t + S_{t^-}(J - 1)\, dP_t.
\). As discussed in Section~\ref{sec:classical_models_literature} the log-return distribution under Merton's model is a Poisson-weighted mixture of
Gaussians. This result requires $\log J$ to be normally distributed; consistent
with Merton's \cite{merton1976option} original formulation, $J$ is therefore \textit{log-normally}
distributed:
\begin{equation}
    \label{eq:jump_lognormal_app}
    \log J \;\sim\; \mathcal{N}(\mu_j,\, \sigma_J^2).
\end{equation}
The parameter $\mu_j$ is the mean log-jump size, while the parameter $\sigma_J^2$ controls the dispersion of jump
sizes.

\subsection{Log-Return Distribution under Merton's Model}

Applying the jump-diffusion analogue of It\^{o}'s Lemma to $f(S_t) = \log S_t$
under Eq.~(\ref{eq:merton_sde_app}), the continuous part is handled exactly as
in Appendix~\ref{app:sec:ito}, while each jump at time $t$ contributes
$\log(J)$ to the log-price. The resulting log-price dynamics are:
\[
    d(\log S_t) = \left(\mu - \frac{1}{2}\nu^2\right)\! dt
                + \nu\, dW_t + \log J\, dP_t.
\]
Integrating over $[0, t]$, conditional on exactly $n$ jumps occurring in
$[0, t]$, the log-return is a sum of the continuous
Gaussian component and $n$ independent log-normal jump contributions:
\[
    \log\!\left(\frac{S_t}{S_0}\right) \;\Bigg|\; P_t = n
    \;\sim\; \mathcal{N}\!\left(\left(\mu - \tfrac{1}{2}\nu^2\right)t
    + n\mu_j,\;\; \nu^2 t + n\sigma_J^2\right).
\]
The unconditional log-return density is obtained by summing over all
possible jump counts, weighted by their Poisson probabilities
$p_n = e^{-\lambda t}(\lambda t)^n / n!$:
\begin{equation}
    \label{eq:merton_return_dist}
    p\!\left(\log\tfrac{S_t}{S_0} = x\right)
    = \sum_{n=0}^{\infty}
    \frac{e^{-\lambda t}(\lambda t)^n}{n!}\;
    \phi\!\left(x;\;
    \left(\mu - \tfrac{1}{2}\nu^2\right)t + n\mu_j,\;
    \nu^2 t + n\sigma_J^2\right),
\end{equation}
where $\phi(\cdot;\, m, v)$ denotes the Gaussian density with mean $m$ and
variance $v$. This is the Poisson-weighted mixture of Gaussians stated in
Section \ref{sec:classical_models_literature}.

\section{Further Limitations of GARCH and Merton's Model}
\label{app:sec:limitations}

Section \ref{sec:classical_models_literature} identifies the core limitation shared by both models, which is further expanded here.

\subsection{Further Limitations of GARCH}

\textbf{Symmetric shock response and the leverage effect} The GARCH(1,1) variance equation $\sigma_t^2 = \omega + \alpha\epsilon_{t-1}^2 + \beta\sigma_{t-1}^2$ depends on $\epsilon_{t-1}^2$, which is symmetric in the sign of the innovation: a negative shock $\epsilon_{t-1} = -c$ and a positive shock $\epsilon_{t-1} = +c$ of the same magnitude produce identical increases in $\sigma_t^2$, contradicting the leverage effect.

\textbf{Tail inadequacy} Even with Student-$t$ innovations, the GARCH conditional tail remains insufficiently heavy. Section \ref{sec:fts_literature} reported power-law tails with exponents $\alpha \in [3,5]$; standard GARCH with Gaussian innovations has finite moments of all orders, and heavier-tailed innovations do not reproduce the empirical tail decay.

\textbf{Short-memory volatility autocorrelation} In GARCH(1,1), squared-return autocorrelation decays exponentially at rate $(\alpha + \beta)^k$. Empirically, volatility autocorrelation decays hyperbolically, $\operatorname{Corr}(|r_t|, |r_{t+k}|) \sim k^{-\beta}$ with $\beta \in [0.1, 0.5]$ \cite{cont2001empirical}, causing underestimation of volatility persistence at long horizons even when $\alpha + \beta$ is close to one.

\textbf{Absence of jumps} GARCH generates smooth conditional variance paths and cannot produce discontinuous price gaps or isolated large moves.

\textbf{Re-estimation requirement} The model must be re-fitted from scratch whenever a new regime arises, with no native mechanism to detect regime transitions.

\subsection{Further Limitations of Merton's Model}

\textbf{Constant jump intensity} The Poisson intensity $\lambda$ in Eq.~(\ref{eq:merton_sde_app}) is fixed, whereas in reality jump activity is time-varying and clustered: large moves arrive in bursts during market stress and are rare during calm periods. A constant $\lambda$ cannot reproduce this self-reinforcing behaviour.

\textbf{Tail inadequacy despite fat tails} Although the Poisson mixture (Eq.~\ref{eq:merton_return_dist}) has heavier tails than a single Gaussian, all moments remain finite. The $n$-th component has variance $\nu^2 t + n\sigma_J^2$ growing linearly in $n$, but the Poisson weights $e^{-\lambda t}(\lambda t)^n/n!$ decay super-exponentially, so the tail of Eq.~(\ref{eq:merton_return_dist}) decays faster than any power law.

\textbf{No volatility clustering} The diffusion component $\nu\, S_t\, dW_t$ has constant volatility $\nu$. Jumps introduce excess kurtosis, but not persistent time-varying conditional variance: between jump events the model reverts to a constant-volatility GBM, and a jump does not alter the volatility of subsequent diffusion increments.

\textbf{No leverage effect} Both the diffusion noise $\nu\, dW_t$ and the jump multiplier $J$ are specified symmetrically with respect to the sign of the return, so Merton's model cannot reproduce the negative lead-lag correlation $L(k) < 0$ defined in Section~\ref{sec:fts_literature}.

\textbf{Calibration difficulty} The likelihood surface of Eq.~(\ref{eq:merton_return_dist}) has multiple local optima, and the jump parameters $(\lambda, \mu_j, \sigma_J^2)$ are poorly identified unless the sample is long and contains clearly distinguishable events. With limited data, a small number of outliers can dominate the estimation of $\lambda$ and $\sigma_J^2$, leading to unstable parameter estimates, which is why we had to constrain the generation process.

\section{Classical Baseline Models: Likelihood and Estimation}
\label{app:baselines}

\subsection{Log-GARCH-X Likelihood}

Given the variance recursion~(\ref{eq:garchx_var}), the standardised residual
is $z_t = \varepsilon_t / \sigma_t$. The negative log-likelihood is:
\begin{equation}
    \label{eq:garchx_nll_app}
    \mathcal{L}_{\text{GARCH-X}} = -\sum_{t=2}^T
        \Bigl[\log p_{t_\nu}(z_t) - \tfrac{1}{2}\log\sigma_t^2\Bigr],
\end{equation}
where $p_{t_\nu}$ is the standardised Student-$t$ log-pdf. The
recursion is initialised at $\sigma_1^2 = \operatorname{Var}(r)$, and the
stability constant $c = 10^{-8}$ prevents $\log(\varepsilon_{t-1}^2 + c)$
from becoming $-\infty$ on zero residuals.

\subsection{Merton-X Likelihood}

Conditional on $P_t = k$ jumps, $r_t \sim \mathcal{N}(\mu + k\mu_J,\,
\nu_t^2 + k\sigma_J^2)$. Marginalising over the Poisson counts:
\begin{equation}
    \label{eq:mertonx_loglik_app}
    \log p(r_t) = \log \sum_{k=0}^{K_{\max}}
        \frac{e^{-\lambda}\lambda^k}{k!}\;
        \mathcal{N}\!\left(r_t;\; \mu + k\mu_J,\;
        \nu_t^2 + k\sigma_J^2\right),
\end{equation}
evaluated via the log-sum-exp identity for numerical stability. The
truncation $K_{\max} = 30$ covers all relevant probability mass for
$\lambda \lesssim 15$; at that intensity, $P(P_t > 30) < 10^{-6}$.
The full-series log-likelihood is a sum of $T$ independent terms (no
sequential dependency).

\subsection{Estimation Procedure}

Both models are fitted per ticker by minimising the negative
log-likelihood with L-BFGS-B \cite{byrd1995limited}. To reduce
sensitivity to initialisation, two structurally distinct starting
points are tried and the one with the lower NLL is retained. For
GARCH-X, $(\alpha_0, \beta_0) \in \{(0.05, 0.90),\, (0.03, 0.90)\}$;
for Merton-X, $(\lambda_0, \sigma_{J,0}) \in \{(0.10, 0.01),\,
(0.50, 0.05)\}$.

\subsection{Implementation Constraints and Design Choices}
\label{app:baselines_constraints}

The following model-specific choices were made to ensure numerical
stability and physical plausibility during fitting and simulation.

\textbf{GARCH-X.} \quad
The model is formulated in log-variance space: the recursion propagates
$\log\sigma_t^2$ rather than $\sigma_t^2$ directly, keeping variance
strictly positive by construction.
Two exogenous regressors enter the variance equation at each time step $t$,
both observable before the day's close: (i)~the squared overnight
log-return $\texttt{open}_t^2$, and (ii)~the Parkinson intra-day variance
estimator $(\log H_t/L_t)^2/(4\log 2)$.
The mean equation includes an AR(1) term whose coefficient $\phi$ is
hard-bounded to $(-0.5,\,0.5)$ to prevent simulated paths from diverging.
Degrees of freedom are constrained to $\nu > 2.01$ via
$\nu = 2.01 + e^{\vartheta}$, guaranteeing finite variance of the
innovation distribution.
Covariance stationarity is enforced by a persistence cap
$\alpha + \beta < 0.98$; any parameter vector violating this bound is
penalised with NLL $= 10^{12}$.
The log-squared lagged residual is regularised as
$\log(\varepsilon_{t-1}^2 + 10^{-8})$ to prevent $\log 0$, and the
log-variance is clipped to $[-30,\, 30]$ to prevent overflow.
The recursion is initialised at $\log\operatorname{Var}(r)$.

\textbf{Merton-X.} \quad
The diffusion variance is time-varying with log-linear dependence on
two intra-day regressors: $\log\nu_t^2 = a_0 + \mathbf{q}_t^\top\mathbf{a}$,
where $\mathbf{q}_t = [\texttt{open}_t^2,\;(\texttt{high}_t -
\texttt{low}_t)^2]^\top$.
Unlike GARCH-X, both regressors are winsorised at their 99th-percentile
training value before entering the likelihood, limiting the influence of
extreme intra-day moves on the fitted variance surface.
Hard bounds are imposed on the jump parameters:
$\log\lambda \in [-6,\,3]$ (i.e.\ $\lambda \in (0.002,\,20)$),
$\mu_J \in [-0.5,\,0.5]$, and $\log\sigma_J \in [-8,\,0]$
(i.e.\ $\sigma_J \leq 1$), preventing degenerate jump configurations
during optimisation.
\textit{Correction to Eq.~\ref{eq:mertonx_loglik_app}}: the Poisson
series is truncated at $K_{\max} = 20$, not $K_{\max} = 30$ as stated
in that equation. For S\&P~500 daily returns, fitted $\lambda$ values
are consistently below $5$, for which $P(N > 20\,|\,\lambda) < 8\times
10^{-8}$; the remaining tail mass is numerically negligible and the
tighter truncation has no effect on results.
